\section{Model Context Protocol (MCP)}

\subsection{Por qué se necesita MCP}

Los LLM tienen un conocimiento del mundo amplio pero \textbf{estático}: solo se actualiza cuando se vuelven a entrenar. Para construir sistemas verdaderamente autónomos se necesita una manera confiable de inyectar datos dinámicos al LLM (clima, precios de acciones, documentación de API más reciente, etc.).

RAG y tool calling son, por ahora, las mejores respuestas. Pero el problema es el siguiente:

\begin{knowledgebox}{El problema del conector $M \times N$}
Supongamos que hay $M$ aplicaciones de LLM y $N$ herramientas/API de terceros. Antes de MCP, cada aplicación tenía que escribir código de integración independiente para cada herramienta, incluida la autenticación, el manejo de errores, la limitación de tasa, etc. Esto generaba una carga de mantenimiento de $M \times N$ conectores.
\end{knowledgebox}

\subsection{Qué es MCP}

MCP (Model Context Protocol) es un protocolo abierto que Anthropic presentó en noviembre de 2024, que ofrece a los LLM un formato \textbf{estandarizado} para exponer herramientas.

\begin{importantbox}{El valor central de MCP}
MCP reduce la integración de herramientas de $M \times N$ a $M + N$:
\begin{itemize}
    \item No hace falta reimplementar la autenticación, el manejo de errores ni la limitación de tasa para cada herramienta
    \item Usa JSON-RPC para forzar un formato de salida consistente
    \item Se extiende a partir del Language Server Protocol (LSP), pero admite un agentic workflow \textbf{proactivo} (proactive), a diferencia del modo puramente reactivo de LSP
\end{itemize}
\end{importantbox}

\subsection{Arquitectura de MCP en detalle}

\subsubsection{Terminología central}

\begin{itemize}
    \item \textbf{Host}: la aplicación anfitriona (como Cursor, Claude Desktop)
    \item \textbf{MCP Client}: la biblioteca integrada en el Host (mantiene una sesión con estado para cada Server)
    \item \textbf{MCP Server}: un envoltorio (wrapper) ligero para la herramienta
    \item \textbf{Tool}: una función que se puede invocar (fuente de datos, API, etc.)
\end{itemize}

\subsubsection{Flujo de trabajo}

\begin{enumerate}
    \item El Client invoca \texttt{tools/list} para preguntarle al MCP Server: «¿qué puedes hacer?»
    \item El Server devuelve una descripción en JSON de cada herramienta (nombre, resumen, JSON Schema)
    \item El Host inyecta ese JSON en el context del modelo
    \item El prompt del usuario activa al modelo, que produce un tool call estructurado
    \item El MCP Server ejecuta la llamada y la conversación continúa
\end{enumerate}

\subsubsection{Capa de transporte}

MCP ofrece dos mecanismos de transporte:
\begin{itemize}
    \item \textbf{stdio}: entrada/salida estándar, adecuada para la comunicación entre procesos locales
    \item \textbf{SSE} (Server-Sent Events): adecuado para la comunicación con servicios remotos
\end{itemize}

\subsection{Limitaciones de MCP}

\begin{warningbox}{La capacidad actual de los Agent para manejar muchas herramientas es limitada}
\begin{itemize}
    \item Los Agent no rinden bien ante una gran cantidad de herramientas: las descripciones de las herramientas consumen rápidamente la context window
    \item El diseño de las API debe orientarse a lo AI-native, en lugar de mantener el diseño rígido de las API tradicionales
    \item Los parámetros de API complejos y las estructuras anidadas reducen de forma notable la precisión de las tool calls
\end{itemize}
\end{warningbox}

\subsection{Resumen de la sección}

MCP es una infraestructura fundamental del ecosistema de herramientas de programación con IA: mediante un protocolo estandarizado resuelve el problema $M \times N$ de la integración entre los LLM y las herramientas externas. Entender la arquitectura de MCP (Host $\rightarrow$ Client $\rightarrow$ Server $\rightarrow$ Tool) y su flujo de comunicación es clave para construir y usar herramientas de programación con IA.
